% © 2026 Ing. David Jaroš — CC BY-NC-ND 4.0
%
% This work is licensed under a Creative Commons Attribution-NonCommercial-NoDerivatives
% 4.0 International License (CC BY-NC-ND 4.0).
%
% Kac-Moody Level k — Approaches Bypassing the Circularity (I2–I5)
% Track: CORE — α derivation — v68 baseline
% Date: 2026-03-09

\ifdefined\INCLUDEMODE
  % Being included in another document — skip preamble
\else
  \documentclass[11pt,a4paper]{article}
  \usepackage{amsmath,amssymb,amsthm}
  \usepackage{hyperref}
  \usepackage{booktabs}

  \newtheorem{theorem}{Theorem}
  \newtheorem{lemma}{Lemma}
  \newtheorem{proposition}{Proposition}
  \newtheorem{conjecture}{Conjecture}
  \theoremstyle{definition}
  \newtheorem{gap}{Gap}
  \theoremstyle{remark}
  \newtheorem{remark}{Remark}
  \newtheorem{observation}{Observation}

  \title{Kac-Moody Level $k$ — Approaches Bypassing the Circularity\\
    \large Approaches I2 (Hosotani Holonomy), I5 (Modular Invariance),\\
    I4 (Anomaly Cancellation), I3 (Representation Theory): v68 Baseline}
  \author{Ing.\ David Jaroš}
  \date{2026-03-09 (v68)}

  \begin{document}
  \maketitle
\fi

%--------------------------------------------------------------------
\section{Problem Statement: The Circularity and How to Bypass It}
\label{sec:intro}
%--------------------------------------------------------------------

\noindent\textbf{No circular reasoning.}
$\alpha^{-1} = 137$ and $n^* = 137$ are \emph{never} used as inputs
anywhere in this document.
$r^2 = \langle\mathrm{Sc}[\Theta^\dagger\Theta]\rangle_{\mathrm{vac}}$
obtained from $V_{\mathrm{eff}}$ is also \emph{not used} (that route
is circular; see below).
Inputs: $\mathcal{S}[\Theta]$, $\theta_H = \pi$ (proved via Hosotani
mechanism), $\psi$-parity (proved symmetry of UBT), $N_{\mathrm{phases}} = 3$
(algebraic fact from $\mathbb{C}\otimes\mathbb{H}$).

\subsection{The fundamental circularity (H-baseline summary)}

The v68 baseline inherits the following established facts and gaps from
\texttt{b\_base\_kac\_moody\_level.tex} (v67):

\begin{itemize}
  \item The formula $k = 2\pi r^2$ was derived exactly from the UBT
    kinetic action. \textbf{[PROVED]}.
  \item The Chern-Simons term is absent from $\mathcal{S}[\Theta]$
    (parity + holonomy arguments). \textbf{[PROVED]}.
  \item $k = 1$ from the minimality principle (free action,
    no CS term). \textbf{[MOTIVATED CONJECTURE]}.
  \item All representation-theory routes (H3a: fundamental, H3b: adjoint,
    H3c: search for $k=1$) yield $k \in \{3/2,\,6\}$ or no match.
    \textbf{[DEAD END for $k=1$]}.
\end{itemize}

The \textbf{fundamental circularity} is the chain:
\[
  k \;\longrightarrow\; c \;\longrightarrow\; B_{\mathrm{base}}
  \;\longrightarrow\; V_{\mathrm{eff}}\;\text{minimum}
  \;\longrightarrow\; n^* \;\longrightarrow\; r^2 \;\longrightarrow\; k.
\]
Any route that determines $k$ from $r^2$ or from $V_{\mathrm{eff}}$ traverses
this circle.  The four approaches investigated here are required to
determine $k$ from the structure of $\mathcal{S}[\Theta]$, the
topology of the $\psi$-circle, or the algebraic properties of the
symmetry algebra — \emph{without} using $r^2$ or $V_{\mathrm{eff}}$
as intermediaries.

\subsection{Approaches investigated (priority order)}

\begin{itemize}
  \item \textbf{I2} (priority 8/10): Topological quantisation +
    Hosotani holonomy $\theta_H = \pi$ (proved) $\to$ winding number
    $w$ $\to$ WZW level $k$.
  \item \textbf{I5} (priority 8/10): Modular invariance of partition
    function + $\psi$-parity symmetry $\to$ uniquely select $k$.
  \item \textbf{I4} (priority 7/10): Anomaly cancellation on the
    $\psi$-circle $\to$ solve for $k$.
  \item \textbf{I3} (priority 7/10): Representation theory of $\Theta$
    under $\mathrm{SU}(2)$ with the Sc[$\cdot$] selector $\to$ $k$.
\end{itemize}

%--------------------------------------------------------------------
\section{Approach I2: Hosotani Holonomy and Topological Winding Number}
\label{sec:I2}
%--------------------------------------------------------------------

\textbf{Status: [DEAD END — winding number $w = 0$; Z$_2$ NS argument
reproduces H2 motivated conjecture but adds no new proof].}

\subsection{Setup: WZW level as a topological invariant}

The WZW level $k$ is a topological invariant for maps
$g : S^3 \to \mathrm{SU}(2) \cong S^3$:
\begin{equation}
  k = \frac{1}{2\pi^2}\int_{S^3}
    \mathrm{Tr}\bigl[(g^{-1}\mathrm{d}g)^3\bigr] \;\in\; \mathbb{Z}.
  \label{eq:k_Pontryagin}
\end{equation}
This is the Pontryagin winding number (degree of the map).
In UBT, $g \in \mathrm{SU}(2)$ is the angular part of the polar
decomposition $\Theta = r\cdot g$ on the $\psi$-circle.
The proved Hosotani result gives:
\begin{equation}
  U \;:=\; \mathrm{P}\exp\!\Bigl(i\oint A_\psi\,\mathrm{d}\psi\Bigr)
  \;=\; \exp\!\bigl(2\pi i\,\theta_H\, T^3\bigr),
  \quad \theta_H = \tfrac{\pi}{2\pi} = \tfrac{1}{2}
  \;\text{(in units of }2\pi),
\end{equation}
i.e.\ $\theta_H = \pi$ in absolute units.

\subsection{Computation: winding number of the Hosotani holonomy}

\begin{observation}[Winding number vanishes for constant holonomy]
  \label{obs:I2_constant}
  The Hosotani vacuum has $A_\psi = \mathrm{const}$, so the holonomy
  $U = \exp(i\pi T^3) \in \mathrm{SU}(2)$ is \emph{spatially constant}
  (independent of 3D position $\mathbf{x}$).
  For a constant $U$, the three-form $\mathrm{Tr}[(U^{-1}\mathrm{d}U)^3]
  = 0$ identically (since $\mathrm{d}U = 0$), and the integral
  in~\eqref{eq:k_Pontryagin} is zero:
  \begin{equation}
    w = \frac{1}{2\pi^2}\int_{S^3}
      \mathrm{Tr}\bigl[(U^{-1}\mathrm{d}U)^3\bigr] = 0.
    \label{eq:w_zero_constant}
  \end{equation}
  Topological non-triviality requires $U(\mathbf{x})$ to vary over 3D
  space and to map $S^3 \to \mathrm{SU}(2)$ with non-zero degree.
  In the constant Hosotani vacuum there is no such map.
\end{observation}

\begin{observation}[Path along $\psi$-circle also gives zero]
  \label{obs:I2_path}
  One might instead consider the map $g(\psi)
  = \exp(i\psi\,\theta_H T^3 / (2\pi))$ for $\psi \in [0,2\pi]$
  as a path in $\mathrm{SU}(2)$.
  Computing the integrand along this path:
  \begin{align}
    g^{-1}\partial_\psi g
    &= \frac{i\theta_H}{2\pi}\, T^3,
    \\
    \mathrm{Tr}\bigl[(g^{-1}\partial_\psi g)^3\bigr]\,\mathrm{d}\psi^3
    &= \Bigl(\frac{i\theta_H}{2\pi}\Bigr)^3\,
      \mathrm{Tr}\bigl[(T^3)^3\bigr]\,\mathrm{d}\psi^3.
  \end{align}
  For $T^3 = \sigma^3/2$ (fundamental): $(T^3)^3 = \sigma^3/8$, and
  $\mathrm{Tr}[\sigma^3/8] = 0$ (the Pauli matrices are traceless).
  The integrand vanishes identically; $w = 0$ regardless of $\theta_H$.

  Physically: $\pi_1(\mathrm{SU}(2)) = \pi_1(S^3) = 0$ — the
  fundamental group of $\mathrm{SU}(2)$ is trivial.  Every loop
  $S^1 \to \mathrm{SU}(2)$ is contractible.  There is no
  integer-valued topological invariant for such maps.
\end{observation}

\subsection{Z\texorpdfstring{$_2$}{2} NS-twist interpretation}

Although the direct winding number is zero, the Hosotani holonomy
$U = \exp(i\pi T^3)$ plays a role as a \emph{Z$_2$ twist} in the
partition function.

In the fundamental representation:
\begin{equation}
  U_{\mathrm{fund}} = \exp\!\bigl(i\tfrac{\pi}{2}\sigma^3\bigr)
  = \mathrm{diag}(i,\,-i),
  \label{eq:U_fund}
\end{equation}
which has $\det = 1$ (hence $U \in \mathrm{SU}(2)$), $\mathrm{Tr} = 0$,
and order 4 ($U^4 = I$).  Under $U^2 = \mathrm{diag}(-1,-1) = -I$,
the field acquires the sign of the $\mathrm{SU}(2)$ center element.

In $\widehat{\mathrm{SU}(2)}_k$ the center element $-I$ acts on the
spin-$j$ representation as $(-1)^{2j}$:
\begin{itemize}
  \item Integer $j$: $(-1)^{2j} = +1$ (Ramond/periodic sector).
  \item Half-integer $j$: $(-1)^{2j} = -1$ (Neveu-Schwarz/anti-periodic
    sector = NS sector).
\end{itemize}
The NS twist $\delta = 1/2$ (proved: Hosotani branch with $\theta_H = \pi$)
projects the Hilbert space onto \emph{half-integer} spin primaries.
For $\widehat{\mathrm{SU}(2)}_k$ the allowed spins are
$j \in \{0, \tfrac{1}{2}, 1, \ldots, \tfrac{k}{2}\}$:
the half-integer NS primaries are $j \in \{\tfrac{1}{2}, \tfrac{3}{2},
\ldots\}$, subject to $j \leq k/2$.

\begin{observation}[NS sector does not uniquely select $k$]
  \label{obs:I2_NS}
  The NS sector exists (contains at least the $j = 1/2$ primary) for
  any $k \geq 1$.  The number of NS primaries is $\lfloor k/2 + 1/2 \rfloor$:
  \begin{center}
  \begin{tabular}{lll}
    \toprule
    $k$ & NS primaries & Count \\
    \midrule
    1 & $j = 1/2$ only & 1 \\
    2 & $j = 1/2$ only & 1 \\
    3 & $j = 1/2,\;3/2$ & 2 \\
    4 & $j = 1/2,\;3/2$ & 2 \\
    5 & $j = 1/2,\;3/2,\;5/2$ & 3 \\
    \bottomrule
  \end{tabular}
  \end{center}
  The Hosotani twist $\theta_H = \pi$ (NS sector) is consistent with any
  $k \geq 1$.  To uniquely select $k = 1$ from the NS sector alone, one
  would need an \emph{external} condition such as
  ``the NS sector contains exactly one primary'' — but this condition
  is not derived from $\mathcal{S}[\Theta]$.
\end{observation}

\begin{remark}[I2 reduces to H2]
  The NS argument gives: $k \geq 1$ (the NS sector is non-empty).
  Combined with the absence of a Chern-Simons term (proved, H2) and the
  free kinetic action, the minimum $k = 1$ is recovered — but this is
  exactly the H2 motivated conjecture.  Approach I2 does not add a new
  independent constraint; it confirms the H2 lower bound $k \geq 1$.
\end{remark}

\textbf{I2 result summary:}
\begin{itemize}
  \item Topological winding number of constant Hosotani holonomy:
    $w = 0$ ($\pi_1(\mathrm{SU}(2)) = 0$).  \textbf{[DEAD END]}.
  \item Path-integral winding along $\psi$-circle: $w = 0$
    (traceless generators). \textbf{[DEAD END]}.
  \item Z$_2$ NS-twist interpretation: consistent with any $k \geq 1$;
    the selection of $k = 1$ requires minimality (= H2).
    \textbf{[MOTIVATED CONJECTURE — same as H2, no new proof]}.
\end{itemize}

%--------------------------------------------------------------------
\section{Approach I5: Modular Invariance and \texorpdfstring{$\psi$}{psi}-Parity}
\label{sec:I5}
%--------------------------------------------------------------------

\textbf{Status: [DEAD END — $\psi$-parity is satisfied by
\emph{all} real modular-invariant partition functions for any $k$].}

\subsection{Setup}

The partition function $Z(\tau)$ of $\widehat{\mathrm{SU}(2)}_k$ on the
torus must be modular invariant.  The $\psi$-parity
$\psi \to -\psi$ is a proved symmetry of the UBT action (the
effective potential $V_{\mathrm{eff}}$ is symmetric under
$\psi \to -\psi$).  On the complex time $\tau = t + i\psi$, the
$\psi$-parity acts as:
\begin{equation}
  \psi \;\to\; -\psi \;:\quad \tau = t + i\psi \;\to\; t - i\psi = \bar\tau
  \;\to\; -\bar\tau
  \quad\text{(with time reversal).}
  \label{eq:psi_parity_action}
\end{equation}
The combined transformation is $\tau \to -\bar\tau$.  The
$\psi$-parity condition on $Z$ is:
\begin{equation}
  Z(\tau) = Z(-\bar\tau).
  \label{eq:psi_parity_Z}
\end{equation}

\subsection{Characters and the map \texorpdfstring{$\tau \to -\bar\tau$}{τ→-τ̄}}

The nome variable is $q = e^{2\pi i \tau}$.  Writing $\tau = \tau_1 + i\tau_2$:
\begin{equation}
  q = e^{-2\pi\tau_2}\,e^{2\pi i\tau_1},
  \qquad
  q_{-\bar\tau} := e^{2\pi i(-\bar\tau)}
  = e^{-2\pi\tau_2}\,e^{-2\pi i\tau_1}
  = q^*,
\end{equation}
so the map $\tau \to -\bar\tau$ corresponds to $q \to q^*$.

The $\widehat{\mathrm{SU}(2)}_k$ characters are:
\begin{equation}
  \chi_j(\tau) = \sum_{r \in \mathbb{Z}} \bigl[
    q^{(r(k+2)+2j+1)^2/(4(k+2))} - q^{(r(k+2)-2j-1)^2/(4(k+2))}
  \bigr]
  \Big/
  \bigl(q^{1/8} - q^{-1/8}\bigr),
\end{equation}
which are convergent series in $q$ with \emph{real} coefficients.
For any $q$ with $|q| < 1$ (i.e.\ $\tau_2 > 0$):
\begin{equation}
  \chi_j(-\bar\tau) = \chi_j(q^*) = \overline{\chi_j(\tau)},
  \label{eq:chi_conjug}
\end{equation}
because the series has real coefficients (complex-conjugation acts only
on the powers $q^n \to (q^*)^n$).

\subsection{All real modular invariants satisfy the \texorpdfstring{$\psi$}{psi}-parity condition}

The modular-invariant partition functions of $\widehat{\mathrm{SU}(2)}_k$
are classified (Cappelli–Itzykson–Zuber A–D–E classification):
\begin{equation}
  Z(\tau) = \sum_{j,\,j'} M_{jj'}\,\chi_j(\tau)\,\overline{\chi_{j'}(\tau)},
  \quad M_{jj'} \in \mathbb{Z}_{\geq 0},\quad M_{00} \geq 1.
  \label{eq:Z_general}
\end{equation}
where the $M$ matrix is symmetric, $M = M^T$ (charge conjugation
invariance).

\begin{proposition}[$\psi$-parity is automatic for all real $Z$]
  \label{prop:psi_parity_automatic}
  For any modular-invariant partition function~\eqref{eq:Z_general}
  with a symmetric $M$-matrix ($M_{jj'} = M_{j'j}$):
  \begin{equation}
    Z(-\bar\tau) = Z(\tau).
  \end{equation}
  \begin{proof}
    Using~\eqref{eq:chi_conjug}:
    \begin{align}
      Z(-\bar\tau)
      &= \sum_{j,j'} M_{jj'}\,\chi_j(-\bar\tau)\,\overline{\chi_{j'}(-\bar\tau)}
       = \sum_{j,j'} M_{jj'}\,\overline{\chi_j(\tau)}\,\chi_{j'}(\tau).
    \end{align}
    Swapping the summation indices $j \leftrightarrow j'$ and using
    $M_{jj'} = M_{j'j}$:
    \begin{equation}
      Z(-\bar\tau)
      = \sum_{j,j'} M_{j'j}\,\overline{\chi_{j'}(\tau)}\,\chi_j(\tau)
      = \sum_{j,j'} M_{jj'}\,\chi_j(\tau)\,\overline{\chi_{j'}(\tau)}
      = Z(\tau). \qedhere
    \end{equation}
  \end{proof}
\end{proposition}

\begin{observation}[$\psi$-parity does not select $k$]
  \label{obs:I5_dead}
  Proposition~\ref{prop:psi_parity_automatic} shows that the condition
  $Z(-\bar\tau) = Z(\tau)$ is satisfied by \emph{every} modular-invariant
  partition function (A-type, D-type, and E-type) for \emph{every}
  level $k \geq 1$.  It therefore provides no constraint that could
  single out a specific value of $k$.
\end{observation}

\begin{remark}[D-type check for $k = 4$]
  The D-type invariant at $k = 4$:
  $Z_D = |\chi_0 + \chi_2|^2 + 2|\chi_1|^2$.
  This has the form~\eqref{eq:Z_general} with
  $M_{00} = M_{02} = M_{20} = M_{22} = M_{11} = 1$ (symmetric), so
  Proposition~\ref{prop:psi_parity_automatic} applies and
  $Z_D(-\bar\tau) = Z_D(\tau)$. ✓
  Similarly for all A-type and E-type invariants.
\end{remark}

\textbf{I5 result summary:}
\begin{itemize}
  \item $\psi$-parity $Z(-\bar\tau) = Z(\tau)$: satisfied by all real
    modular-invariant partition functions for any $k$.
    \textbf{[DEAD END — no constraint on $k$]}.
\end{itemize}

%--------------------------------------------------------------------
\section{Approach I4: Anomaly Cancellation on the \texorpdfstring{$\psi$}{psi}-Circle}
\label{sec:I4}
%--------------------------------------------------------------------

\textbf{Status: [DEAD END — non-chiral bosonic field has no 2D gauge
anomaly; matter-content counting reduces to H3].}

\subsection{Setup: 2D chiral anomaly on $S^1_\psi$}

The WZW level $k$ is related to the 2D gauge anomaly:
in a chiral (2D) gauge theory, the anomaly coefficient equals the
level $k$ of the surviving current algebra.  The condition of anomaly
cancellation reads:
\begin{equation}
  k_{\mathrm{total}} = \sum_{\text{left}} n_i\,T(\mathbf{r}_i)
    - \sum_{\text{right}} n_i\,T(\mathbf{r}_i) = 0,
  \label{eq:anomaly_cancel}
\end{equation}
where left/right refer to the chirality of the mode on $S^1_\psi$,
$n_i$ is the multiplicity, and $T(\mathbf{r}_i)$ is the Dynkin index
of representation $\mathbf{r}_i$.

\subsection{UBT field \texorpdfstring{$\Theta$}{Theta} is non-chiral}

The UBT field $\Theta \in \mathbb{C}\otimes\mathbb{H}$ is a
\emph{complex bosonic} field.  Its kinetic action
$\mathcal{S}_{\mathrm{kin}} = \int \mathrm{Sc}[(D_\mu\Theta)^\dagger(D^\mu\Theta)]\,\mathrm{d}^5x$
is:
\begin{enumerate}
  \item \textbf{Bosonic}: no intrinsic chirality from spin statistics.
  \item \textbf{Real} (under the biquaternion norm): $\mathcal{S}_{\mathrm{kin}}
    = \frac{1}{2}\int \mathrm{Tr}_2[(D\Theta)^\dagger(D\Theta)]\,\mathrm{d}^5x$,
    which treats left- and right-moving modes on $S^1_\psi$ symmetrically.
  \item \textbf{Non-chiral}: for each KK mode $\Theta_n$ with
    $\psi$-momentum $p_\psi = n/R$ on $S^1_\psi$, the complex conjugate
    mode $\Theta_{-n}$ carries momentum $-n/R$.
    The left-mover (momentum $+|n|/R$) and right-mover (momentum $-|n|/R$)
    appear in equal numbers with the same $\mathrm{SU}(2)$ representation.
\end{enumerate}

\begin{proposition}[Anomaly vanishes for $\Theta$]
  \label{prop:no_anomaly}
  The 2D $\mathrm{SU}(2)$ gauge anomaly on $S^1_\psi$ vanishes
  identically for the UBT field $\Theta$:
  \begin{equation}
    k_{\mathrm{anom}} =
    \sum_{\text{left}} n_i\,T(\mathbf{r}_i)
    - \sum_{\text{right}} n_i\,T(\mathbf{r}_i) = 0.
  \end{equation}
  \begin{proof}
    Each KK mode $\Theta_n$ with representation $\mathbf{r}$ and
    left-mover contribution $+T(\mathbf{r})$ is accompanied by its
    conjugate $\Theta_{-n}$ with equal representation $\mathbf{r}$ and
    right-mover contribution $-T(\mathbf{r})$.  The contributions cancel
    pairwise.  The $n = 0$ (zero-mode) is non-chiral by construction.
    Hence $k_{\mathrm{anom}} = 0$ for all representations.
  \end{proof}
\end{proposition}

\begin{remark}[NS shift does not restore chirality]
  The Hosotani NS shift $\delta = 1/2$ changes the KK momenta from
  $n/R$ to $(n + 1/2)/R$ (half-integer spectrum).  This does not
  introduce chirality: the shifted mode $(n+1/2)/R$ still has an
  equal-weight partner $-(n+1/2)/R = -(n+1+1/2)/R$ after the shift.
  Anomaly cancellation remains trivial.
\end{remark}

\subsection{Matter-content counting: reduces to H3}

If one ignores chirality and instead asks ``what WZW level is
generated by the matter content on $S^1_\psi$?'', the answer is
determined purely by the Dynkin index sum:
\begin{equation}
  k = \sum_i n_i\,T(\mathbf{r}_i),
  \label{eq:k_dynkin_I4}
\end{equation}
which is \emph{exactly} the H3 formula.  The results from H3 apply
directly:
\begin{itemize}
  \item $N_{\mathrm{phases}} = 3$ modes in the adjoint ($T = 2$):
    $k = 3 \times 2 = 6$.  $B \approx 93.5 \neq 41.57$.
    \textbf{[DEAD END]}.
  \item $N_{\mathrm{phases}} = 3$ modes in the fundamental ($T = 1/2$):
    $k = 3 \times 1/2 = 3/2$.  $B \approx 53.4 \neq 41.57$.
    \textbf{[DEAD END]}.
\end{itemize}
No anomaly-cancellation argument changes these counts.

\textbf{I4 result summary:}
\begin{itemize}
  \item 2D gauge anomaly of the UBT field $\Theta$: vanishes identically
    ($k_{\mathrm{anom}} = 0$) for non-chiral bosons.
    \textbf{[DEAD END — trivial]}.
  \item Matter-content counting: identical to H3 (adjoint gives $k=6$;
    fundamental gives $k=3/2$; neither is $k=1$).
    \textbf{[DEAD END — same as H3]}.
\end{itemize}

%--------------------------------------------------------------------
\section{Approach I3: Representation Theory of \texorpdfstring{$\Theta$}{Theta}
with the Sc\texorpdfstring{$[\cdot]$}{[.]} Selector}
\label{sec:I3}
%--------------------------------------------------------------------

\textbf{Status: [DEAD END — Sc$[\cdot]$ does not change the level;
all natural assignments give $k \in \{2, 6, 3/2\}$, none gives
$k = 1$].}

\subsection{Role of the scalar selector \texorpdfstring{Sc$[\cdot]$}{Sc[.]}}

The UBT kinetic action uses the \emph{scalar part selector}:
\begin{equation}
  \mathrm{Sc}[Q] = \tfrac{1}{2}\mathrm{Tr}_2[Q]
  \quad \text{for } Q \in \mathrm{Mat}(2,\mathbb{C}),
\end{equation}
where $\mathrm{Tr}_2$ is the $2\times 2$ matrix trace.
Thus:
\begin{equation}
  \mathrm{Sc}\bigl[(D\Theta)^\dagger(D\Theta)\bigr]
  = \tfrac{1}{2}\,\mathrm{Tr}_2\bigl[(D\Theta)^\dagger(D\Theta)\bigr]
  = \tfrac{1}{2}\sum_{ij}\bigl|(D\Theta)_{ij}\bigr|^2.
  \label{eq:Sc_as_trace}
\end{equation}
This is the Hilbert–Schmidt (Frobenius) norm squared of $D\Theta$,
summed over all $2\times 2$ matrix entries.  The Sc$[\cdot]$ selector
does \emph{not} project onto a group-theoretic singlet; it is the
standard $\mathrm{SU}(2)$-invariant trace norm.

\subsection{SU(2) representations of \texorpdfstring{$\Theta$}{Theta}}

$\Theta \in \mathbb{C}\otimes\mathbb{H} \cong \mathrm{Mat}(2,\mathbb{C})$
has $2\times 2 = 4$ complex components.

\paragraph{I3a: Left-action $g \cdot \Theta = g\Theta$ (SU(2)$_L$).}
Under $g \to g\Theta$, each \emph{column} of the $2\times 2$ matrix
$\Theta$ transforms as the fundamental (spin-$\tfrac{1}{2}$)
representation of $\mathrm{SU}(2)_L$.  There are 2 complex columns,
each containing 2 complex components:
\begin{itemize}
  \item Two complex fundamentals: $n_{\mathrm{fund,}\mathbb{C}} = 2$.
  \item Per complex fundamental: contribution $T(\mathrm{fund})_\mathbb{C}
    = 2\,T(\mathrm{fund})_\mathbb{R} = 2\times\tfrac{1}{2} = 1$.
  \item Total: $k_{I3a} = 2 \times 1 = 2$.
\end{itemize}
$c_{k=2} = 3\times 2/(2+2) = 3/2$.
$B_{\mathrm{base}} = (3/2)\times 12^{3/2} \approx 62.4 \neq 41.57$.
\textbf{[DEAD END]}.

\paragraph{I3b: Adjoint action $g \cdot \Theta \cdot g^{-1}$
(SU(2)$_{\mathrm{diag}}$).}
The decomposition $\mathrm{Mat}(2,\mathbb{C}) = \mathbb{C}\otimes\mathbb{H}
= \mathrm{span}\{1\} \oplus \mathrm{span}\{\mathbf{i},\mathbf{j},\mathbf{k}\}
= \mathrm{singlet} \oplus \mathrm{adjoint}$
under the diagonal $\mathrm{SU}(2)$.
The imaginary modes $\mathrm{Im}(\mathbb{H}) = \mathrm{span}\{\mathbf{i},
\mathbf{j},\mathbf{k}\}$ transform as the adjoint (spin-1) with $T(\mathbf{adj}) = 2$.
\begin{itemize}
  \item $N_{\mathrm{phases}} = 3$ real adjoint scalars.
  \item $k_{I3b} = 3 \times T(\mathbf{adj}) = 3 \times 2 = 6$.
    (Same as H3b.)
\end{itemize}
$c_{k=6} = 3\times 6/8 = 9/4$.
$B_{\mathrm{base}} = (9/4)\times 12^{3/2} \approx 93.5 \neq 41.57$.
\textbf{[DEAD END — same as H3b]}.

\paragraph{I3c: Can Sc$[\cdot]$ reduce the effective level?}
One might hope that since Sc$[D\Theta^\dagger D\Theta]
= \frac{1}{2}\mathrm{Tr}_2[(D\Theta)^\dagger(D\Theta)]$
selects only the \emph{trace} of the matrix product, the WZW
current might be effectively smaller.

The $\mathrm{SU}(2)_L$-current on $S^1_\psi$ is:
\begin{equation}
  J^a(\psi)
  = \frac{\delta \mathcal{S}_{\mathrm{kin}}}{\delta A^a_\psi}
  = \mathrm{Sc}[\Theta^\dagger T^a \partial_\psi \Theta]
    + \mathrm{h.c.}
  = \tfrac{1}{2}\,\mathrm{Tr}_2[\Theta^\dagger T^a\partial_\psi\Theta]
    + \mathrm{h.c.}
  \label{eq:J_current}
\end{equation}
This is exactly the standard $\mathrm{SU}(2)_L$-current for
$\Theta$ in the 2-column-fundamental representation.
The OPE $J^a(\psi)J^b(\psi') \sim k\,\delta^{ab}/(2\pi(\psi-\psi')^2)$
gives $k = k_{I3a} = 2$ (from the two complex fundamental columns),
or the adjoint counting $k = 6$, depending on which
$\mathrm{SU}(2)$ action is gauged.

The Sc$[\cdot]$ selector does not reduce $k$:
it is the standard trace that makes the action $\mathrm{SU}(2)$-invariant,
not a projection that removes degrees of freedom.

\begin{observation}[Sc$[\cdot]$ cannot give $k=1$]
  \label{obs:I3_no_k1}
  The four natural assignments:
  \begin{center}
  \begin{tabular}{lllll}
    \toprule
    Label & $\mathrm{SU}(2)$ action & Rep & $k$ & $B_{\mathrm{base}}$ \\
    \midrule
    I3a & left: $g\Theta$ & 2 complex fund.\ & 2 & 62.4 \\
    I3b & adjoint: $g\Theta g^{-1}$ & 3 real adj.\ & 6 & 93.5 \\
    H3a & adjoint on $\mathrm{Im}(\mathbb{H})$ & 3 real fund.\ & 3/2 & 53.4 \\
    H3b & adjoint on $\mathrm{Im}(\mathbb{H})$ & 3 real adj.\ & 6 & 93.5 \\
    \bottomrule
  \end{tabular}
  \end{center}
  None of these gives $k = 1$.
  For $k = 1$ one would need $\sum_i n_i T(\mathbf{r}_i) = 1$;
  the only option with small integers is $n = 2$, $T = 1/2$
  (two real fundamentals), but $\mathbb{C}\otimes\mathbb{H}$ does not
  naturally provide exactly 2 real fundamental modes (it has 3 imaginary
  + 1 scalar = 4 quaternionic components).
\end{observation}

\textbf{I3 result summary:}
\begin{itemize}
  \item Left-action (I3a): $k = 2$; $B \approx 62.4 \neq 41.57$.
    \textbf{[DEAD END]}.
  \item Adjoint-action (I3b): $k = 6$; $B \approx 93.5 \neq 41.57$.
    (Identical to H3b.) \textbf{[DEAD END]}.
  \item Sc$[\cdot]$ selector: does not reduce $k$; it is the standard
    SU(2)-invariant trace norm.  \textbf{[DEAD END]}.
  \item Search for $k = 1$: no natural assignment from
    $\mathbb{C}\otimes\mathbb{H}$ achieves $\sum n_i T_i = 1$.
    \textbf{[DEAD END]}.
\end{itemize}

%--------------------------------------------------------------------
\section{Summary and Cumulative Status A1--I5}
\label{sec:summary}
%--------------------------------------------------------------------

\subsection{Approaches I2--I5: results table}

\begin{center}
\begin{tabular}{p{1.4cm}p{4.2cm}p{4.5cm}p{3.2cm}}
  \toprule
  Approach & Method & Conclusion & Status \\
  \midrule
  I2 (topo.) & Pontryagin winding of $U_{\mathrm{Hosotani}}$ &
    $w = 0$ ($\pi_1(\mathrm{SU}(2)) = 0$; constant holonomy) &
    \textbf{[DEAD END]} \\[4pt]
  I2 (Z$_2$) & NS twist: $\theta_H = \pi$ selects half-int.\ sector &
    $k \geq 1$ consistent; same as H2 minimality &
    \textbf{[MOTIVATED CONJECTURE (= H2)]} \\[4pt]
  I5 & $\psi$-parity $Z(-\bar\tau) = Z(\tau)$ &
    Satisfied by all real invariants for any $k$ &
    \textbf{[DEAD END]} \\[4pt]
  I4 (anom.) & Chiral anomaly on $S^1_\psi$ &
    Vanishes identically for non-chiral boson $\Theta$ &
    \textbf{[DEAD END — trivial]} \\[4pt]
  I4 (matter) & Dynkin-index counting (as in H3) &
    $k = 6$ (adj.) or $k = 3/2$ (fund.); neither $= 1$ &
    \textbf{[DEAD END (= H3)]} \\[4pt]
  I3a & Left action $g\Theta$: 2 complex fund.\ &
    $k = 2$; $B \approx 62.4 \neq 41.57$ &
    \textbf{[DEAD END]} \\[4pt]
  I3b & Adjoint $g\Theta g^{-1}$: 3 real adj.\ &
    $k = 6$; $B \approx 93.5$ (= H3b) &
    \textbf{[DEAD END (= H3b)]} \\[4pt]
  I3c & Sc$[\cdot]$ selector: role in level &
    Sc is standard trace; does not reduce $k$ &
    \textbf{[DEAD END]} \\
  \bottomrule
\end{tabular}
\end{center}

\subsection{Cumulative status A1--I5 (30 approaches)}

\begin{center}
\begin{tabular}{p{2.0cm}p{5.0cm}p{5.5cm}}
  \toprule
  Approach group & Method family & Best status \\
  \midrule
  A1--F4 (v65) & Gaussian PI, Weyl anom., Fueter, QK index, holonomy,
    modular weight, \ldots & \textbf{[DEAD END]} or \textbf{[NUMERICAL OBS]} \\
  G3 (v65) & Kac-Moody modular anomaly & \textbf{[PARTIAL]: $k = 1$ consistent, $c = 1$} \\
  G2, G4, G7, F (v65) & Weyl $\tilde c$, Fueter, QK index, non-isotropic $D$ &
    \textbf{[DEAD END]} \\
  H1 (v67) & WZW norm.: $k = 2\pi r^2$, $r^2 = f(n^*)$ &
    \textbf{[DEAD END (circular)]} \\
  H2 (v67) & No CS + free action + minimality &
    \textbf{[MOTIVATED CONJECTURE]: $k = 1$} \\
  H3 (v67) & Dynkin index / rep.\ theory &
    \textbf{[DEAD END]: $k \in \{3/2,\,6\}$} \\
  I2 (v68) & Hosotani holonomy winding &
    \textbf{[DEAD END]} ($w = 0$); Z$_2$ reduces to H2 \\
  I5 (v68) & Modular inv.\ + $\psi$-parity &
    \textbf{[DEAD END]}: all real $Z$ satisfy cond.\ \\
  I4 (v68) & Anomaly cancellation &
    \textbf{[DEAD END]}: trivial for non-chiral boson \\
  I3 (v68) & Rep.\ theory + Sc$[\cdot]$ selector &
    \textbf{[DEAD END]}: $k \in \{2,\,6\}$; none $= 1$ \\
  \midrule
  \multicolumn{2}{l}{\textbf{Overall best status (30 approaches total):}} &
    $k = 1$ is \textbf{[MOTIVATED CONJECTURE]} \\
  \multicolumn{2}{l}{\textbf{Gap (G3-k) status:}} &
    \textbf{[OPEN — see updated gap below]} \\
  \bottomrule
\end{tabular}
\end{center}

\subsection{Updated gap status after I2--I5}

\begin{gap}[Gap (G3-k): Level of $\widehat{\mathrm{SU}(2)}$ in UBT —
Updated after Approaches I2--I5]
  \label{gap:G3k_I5}
  After 30 approaches (A1--I5), the situation is unchanged from v67:
  the formula $k = 2\pi r^2$ is proved; the CS term is absent;
  $k = 1$ is a motivated conjecture from the minimality principle (H2).
  
  Approaches I2--I5 have each been carried through to a conclusion:
  \begin{itemize}
    \item I2 (topological winding): $w = 0$ because $\pi_1(\mathrm{SU}(2)) = 0$.
      The NS-sector interpretation reproduces H2 and does not add a new proof.
    \item I5 (modular invariance + $\psi$-parity): the ψ-parity condition
      is satisfied by all real modular-invariant partition functions,
      providing no constraint on $k$.
    \item I4 (anomaly cancellation): the 2D anomaly is trivially zero for
      the non-chiral bosonic field $\Theta$; the matter-content counting
      reproduces the H3 dead ends.
    \item I3 (representation theory with Sc$[\cdot]$): the left action
      on $\mathrm{Mat}(2,\mathbb{C})$ gives $k = 2$; the adjoint gives
      $k = 6$; the Sc$[\cdot]$ selector is the standard trace and does
      not reduce $k$.  No natural assignment gives $k = 1$.
  \end{itemize}
  \textbf{All four new routes are [DEAD END].}
  
  The only path to $k = 1$ that is not a dead end remains H2
  (minimality: no CS, free action, $k = 1$ as the lowest consistent
  level).  This is a \textbf{[MOTIVATED CONJECTURE]}, not a proof.
  
  The fundamental obstruction is: any non-circular derivation of $k$
  must either (a)~determine $r^2$ from $\mathcal{S}[\Theta]$ without
  $V_{\mathrm{eff}}$, or (b)~find a topological or algebraic invariant
  of $\mathcal{S}[\Theta]$ that is integer-valued and equals 1.
  Route (a) remains circular; route (b) has not been found despite
  30 attempts.
  
  Status: \textbf{[OPEN — $k=1$ remains [MOTIVATED CONJECTURE];
  new approaches needed]}.
\end{gap}

\subsection{What has been gained from I2--I5}

The negative results of I2--I5 are informative:
\begin{enumerate}
  \item The simplest topological invariant of the Hosotani configuration
    ($w = 0$) is trivial, ruling out the topology of $S^1_\psi$ as the
    source of the quantisation of $k$.
  \item The ψ-parity condition, while a genuine constraint on the
    UBT action, does not translate into a constraint on $k$ — it is
    automatically satisfied by the symmetry of ADE modular invariants.
  \item The 2D gauge anomaly on $S^1_\psi$ is trivial for bosons;
    chirality (which would create a non-trivial anomaly) is absent.
  \item The scalar-part selector Sc$[\cdot]$ in the UBT action is
    equivalent to the standard $\mathrm{SU}(2)$-invariant trace and
    does not project the level to a lower value.
\end{enumerate}
Taken together, these results sharpen the obstruction: $k$ is not
determined by the topology, symmetry, anomaly structure, or
representation content of the free kinetic action alone.
The only remaining handle is the \emph{normalisation} of the kinetic
term ($k = 2\pi r^2$), which requires an independent source for
$r^2$.

%--------------------------------------------------------------------
\section{In Plain Language}
\label{sec:plain}
%--------------------------------------------------------------------

\textbf{In plain language:}
This document investigated four new approaches to determine the integer
$k$ (the ``Kac-Moody level'' of the imaginary-time symmetry algebra of
UBT), all designed to avoid the circularity that blocks the earlier
approaches.

The first new approach (I2) tried to use the known value of the
gauge-field phase around the imaginary-time circle (called the
Hosotani angle, proved to be $\pi$) to extract a ``winding number''
that would equal $k$.
The attempt fails because the gauge group $\mathrm{SU}(2)$ is
topologically a 3-sphere, which has no non-trivial loops:
every path in $\mathrm{SU}(2)$ that starts and ends at the same point
can be continuously shrunk to a point.
Mathematically, $\pi_1(\mathrm{SU}(2)) = 0$.
The winding number of the Hosotani holonomy is therefore zero,
not $k = 1$ or any other positive integer.
A secondary argument — that the Hosotani angle $\pi$ activates the
``anti-periodic'' (NS) sector of the algebra, which exists only for
$k \geq 1$ — correctly gives a lower bound $k \geq 1$, but this is
exactly the same result already obtained by the minimality argument
(Approach H2).

The second new approach (I5) tested whether the symmetry of the
partition function under time-reversal (the ``$\psi$-parity'' symmetry
proved for UBT) could single out a unique level $k$.
The answer is no: every physically sensible partition function
(technically: every modular-invariant $Z$ with a symmetric coupling
matrix) automatically satisfies the time-reversal condition.
This is a mathematical identity that holds for every value of $k$
simultaneously, so no value is preferred.

The third approach (I4) asked whether requiring the quantum theory on
the imaginary-time circle to be free of gauge anomalies fixes $k$.
For a bosonic field like the UBT field $\Theta$, the gauge anomaly
is automatically zero (left- and right-moving modes come in pairs
of equal and opposite charge, cancelling each other).
When the anomaly is non-trivially extended to a level-counting
formula, the answer reduces to the same mode-counting that was already
exhausted in the representation-theory approach H3.

The fourth approach (I3) revisited the representation theory of the
UBT field under $\mathrm{SU}(2)$, paying special attention to the
role of the scalar-part selector Sc$[\cdot]$ in the action.
We found that Sc$[\cdot]$ is algebraically equivalent to the standard
matrix trace, which is the $\mathrm{SU}(2)$-invariant inner product.
It does not reduce the effective level: the left-multiplication
action gives $k = 2$, and the adjoint action gives $k = 6$,
neither matching the needed $k = 1$.

After 30 approaches, $k = 1$ remains a well-motivated but unproved
conjecture.
The precise obstacle is now clearly identified: all routes to $k$
either require knowing the vacuum amplitude $r^2$
(which cannot be obtained without the very number $n^* = 137$
that we are trying to derive), or produce integer values different
from 1.
A proof of $k = 1$ will require either a novel algebraic structure
in $\mathbb{C}\otimes\mathbb{H}$ that is not captured by the standard
representation theory, or an entirely new strategy.

\ifdefined\INCLUDEMODE\else
\end{document}
\fi
